\chapter{Wstęp i cel pracy}
\label{chap:introduction}

Celem projektu jest stworzenie gry geograficzno-geolokalizacyjnej polegającej na odgadywaniu miejsc związanych z Politechniką Gdańską na podstawie ich zdjęć. Dodatkowo dostępny będzie tryb terenowy, w którym użytkownik musi fizycznie udać się w odgadywane miejsce. Aby zwiększyć stopień interakcji między użytkownikami, zaplanowaliśmy również udostępnienie publicznego rankingu, zachęcającego do poprawiania swoich wyników, a także dodanie trybu wieloosobowego, w którym gracze rywalizowaliby między sobą, ścigając się o to, kto pierwszy odgadnie zadaną lokalizację.

Aplikacja taka pozwalałaby na interaktywne poznawanie terenów kampusu dla szerokiego grona zainteresowanych, od osób niezwiązanych z uczelnią, do absolwentów oraz pracowników PG. Szczególnie przydatnym zastosowaniem byłoby zapoznawanie przyszłych i nowych studentów z okolicami uczelni, ułatwiając im orientację w nowym otoczeniu.

Gry przeglądarkowe cieszyły się popularnością od momentu ich powstania. Ich głównym atutem są niskie wymagania sprzętowe czy brak potrzeby instalacji, które od początku poszerzyły grono potencjalnych odbiorców. Dodatkowo, użytkownicy za pośrednictwem dostępu do internetu mogą wchodzić ze sobą w interakcje i rywalizować dzięki udostępnianym rankingom.

Popularność gier komputerowych znacząco wzrosła podczas pandemii COVID-19, gdzie przez potrzebę pozostania w izolacji możliwości spędzania wolnego czasu były ograniczone. Wydłużony czas pozostania w domach przyczynił się też do wzrostu zainteresowania grami geograficznymi i geolokalizacyjnymi. Jednym z czołowych produktów tego zjawiska został GeoGuessr, który mimo wypuszczenia na rynek 6 lat wcześniej, został spopularyzowany dopiero w 2020 roku. Sukces ten można przypisać idei połączenia rozrywki i odkrywania świata z wykorzystaniem Google Street View, dzięki której użytkownicy mogą się wykazać wiedzą na temat architektury, ukształtowania terenu oraz panującego w danym kraju języka.

Podobnym trendem wykazały się aplikacje mobilne, które ze względu na szybki rozwój urządzeń przenośnych, pozwalają graczowi połączyć się z dowolnego miejsca, w dowolnej chwili. Zlikwidowały one potrzebę posiadania drogiego i nieporęcznego sprzętu, stawiając na wygodę użytkownika. Atuty te wykorzystała gra Pokemon Go, która z pomocą dobrze znanej marki podbiła rynek, osiągając pułap 500 milionów pobrań w pierwszym roku po debiucie. Gra ta korzysta z funkcji GPS do znalezienia wirtualnych aren rozsianych po prawdziwych lokalizacjach, przyczyniając się do zwiększenia immersji rozgrywki.

Po dokonaniu powyższej analizy rynku, postanowiliśmy w ramach projektu inżynierskiego stworzyć aplikację łączącą w sobie dotychczas wymienione aspekty.
